\chapter{The model architecture}
\label{ch:model}

The recurrent vision transformer has four stages: a patch front-end, a recurrent
self-attention block, a spatial recurrent memory, and an actor--critic readout. This
chapter gives the canonical (paper) form of each; later chapters vary one stage at a time.

\section{Patch front-end}
Each $50\times50$ frame is divided into four $25\times25$ quadrants. In the canonical
network each quadrant is encoded by the encoder of a variational autoencoder (VAE) up to
its second flattened layer, a $128$-dimensional feature $\oflat$ (the paper feeds the ViT
this layer, not the stochastic latent). Per-patch positional and per-frame temporal
one-hot codes are concatenated, giving a $140$-dimensional token, and the four tokens form
$\X^{(t)}\in\mathbb R^{4\times140}$. The role of the VAE --- and the consequence of
training it from reward versus from reconstruction --- is the subject of
Chapter~\ref{ch:perception}.

\section{Recurrent multiplicative self-attention}
The canonical attention is single-head and \emph{multiplicatively gated} by memory. With
the previous-frame memory $\Hone^{(t-1)}\in\mathbb R^{4\times\dmem}$ ($\dmem=1024$),
\begin{align}
\mathbf Q &= (\X\,\mathbf W_{XQ})\odot(\Hone\,\mathbf W_{HQ}), &
\mathbf K &= (\X\,\mathbf W_{XK})\odot(\Hone\,\mathbf W_{HK}), &
\mathbf V &= (\X\,\mathbf W_{XV})\odot(\Hone\,\mathbf W_{HV}),
\label{eq:mult}
\end{align}
with $\mathbf W_{X\cdot}\in\mathbb R^{140\times140}$, $\mathbf W_{H\cdot}\in\mathbb
R^{1024\times140}$, $\odot$ the Hadamard product, and the attention
\begin{equation}
\mathbf V_{\text{filt}}=\mathrm{softmax}(\mathbf Q\mathbf K^\top)\mathbf V,\qquad
\mathbf Z^{(t)}=\X^{(t)}+\mathbf V_{\text{filt}}\in\mathbb R^{4\times140}.
\label{eq:z}
\end{equation}
The multiplicative form is the program's neuro-motivated gain-modulation: memory does not
add to the image, it \emph{scales} the bottom-up query/key/value projections. As
Chapter~\ref{ch:attention} documents, this exact form has an important consequence ---
when memory is uninformative the gate is near zero and the attention is uniform --- which
is at the heart of the attention/memory dissociation. The cross-attention and FiLM
variants replace Eq.~\eqref{eq:mult} with, respectively, $\mathbf K,\mathbf V$ formed by
concatenating image and memory tokens, and a $(1+\Hone\mathbf W_{H})$ identity-floored
gate.

\section{Spatial recurrent memory}
The percept $\mathbf Z$ writes a per-patch memory through an exponential-gated spatial
LSTM (an adaptation of the xLSTM), maintained independently per patch so that
self-attention is the only route by which spatial locations exchange information:
\begin{align*}
\tilde{\mathbf I}&=\mathbf Z\mathbf W_i+\mathbf H\mathbf R_i, &
\tilde{\mathbf F}&=\mathbf Z\mathbf W_f+\mathbf H\mathbf R_f, &
\tilde{\mathbf O}&=\mathbf Z\mathbf W_o+\mathbf H\mathbf R_o, &
\tilde{\mathbf U}&=\mathbf Z\mathbf W_u+\mathbf H\mathbf R_z,\\
\mathbf M&=\max(\tilde{\mathbf F}+\mathbf M',\tilde{\mathbf I}), &
\mathbf I&=e^{\tilde{\mathbf I}-\mathbf M}, &
\mathbf F&=e^{\tilde{\mathbf F}+\mathbf M'-\mathbf M}, &
\mathbf O&=\sigma(\tilde{\mathbf O}),
\end{align*}
\[
\mathbf N=\mathbf F\odot\mathbf N'+\mathbf I,\quad
\mathbf C=\mathbf C'\odot\mathbf F+\tanh(\tilde{\mathbf U})\odot\mathbf I,\quad
\mathbf H=\mathbf O\odot(\mathbf C/\mathbf N),
\]
with $\dmem=1024$, $\mathbf W_x\in\mathbb R^{140\times1024}$, $\mathbf
R_x\in\mathbb R^{1024\times1024}$, and the stabiliser/normaliser states $\mathbf M,\mathbf
N$. The memory is the agent's working memory: it carries the cue and the accumulating
evidence across the blank interval.

\section{Actor--critic readout}
The four-patch structure is retained through attention and memory but \emph{flattened} at
the reinforcement-learning agent: the readout is $\mathbf H'\in\mathbb R^{4096}$. A
four-layer feed-forward actor maps $\mathbf H'$ to a two-way policy (wait / declare); a
distributional critic maps $\mathbf H'$ to a value distribution. Which internal state is
read --- the flattened memory $\mathbf H$ (canonical) versus the attention output
$\mathbf Z$ (the cross-talk and codebook variants) --- is itself an architectural variable
that, we find, changes the phenotype.

\begin{table}[t]
\centering
\caption{Components of the canonical recurrent vision transformer and the axes along which
the architecture zoo varies them.}
\label{tab:components}
\small
\begin{tabular}{p{2.6cm}p{4.6cm}p{5.4cm}}
\toprule
Stage & Canonical form & Variation axis explored \\
\midrule
Front-end & VAE encoder $\to\oflat$ (128-d), 4 tokens & raw-pixel vs.\ conv vs.\ frozen
VAE; 4 vs.\ 100 tokens \\
Attention & multiplicative gate, single head & cross-attention, FiLM $(1{+}\cdot)$,
fast-weights, codebook, dual-head \\
Feedback source & memory $\Hone$ scales Q/K/V & memory as keys/values; memory generates
weights; memory as queries \\
Memory & one xLSTM, $\dmem{=}1024$ & two LSTMs (feedback vs.\ decision); parallel vs.\
stacked; LSTM vs.\ vanilla \\
Readout & flattened $\mathbf H$ (4096) & read $\mathbf Z$; split $\mu/q$; codebook
selection \\
\bottomrule
\end{tabular}
\end{table}
